% Source: https://artofproblemsolving.com/community/c3h3260311_mathcounts_cutoff
% Stable identity: aops-topic-3260311-q10
The hypotenuse is $2\sqrt2$, so the circumradius is $R=\sqrt2$. The inradius is $r=\dfrac{2+2-2\sqrt2}{2}=2-\sqrt2$. Thus the difference is $\pi(R^2-r^2)=\boxed{4\pi(\sqrt2-1)}$.
\begin{asy}
usepackage("amsmath");
size(7cm);
pair A=(0,0),B=(2,0),C=(0,2),O=(1,1);real r=2-sqrt(2);draw(circle(O,sqrt(2)),gray);draw(A--B--C--cycle);draw(circle((r,r),r),blue);label("$2$",A--B,S);label("$2$",A--C,W);label("$O$",O,N);
\end{asy}
